\scanpage{85}
\deleted{\unsure{sont les suivantes}}\inserted{s'expriment} comme suit : Si $\sigma\in\mathfrak S_3$, et $\widetilde\sigma_\Omega$ est l'op. sur $\Omega_0$ (ou $\Omega$), on a

\[
 \left\{\begin{aligned}
 \widetilde\sigma_\Omega\,\rho^I(\varepsilon)\,\sigma_\Omega^{-1}
 &=\rho^{\sigma I}(\sigma\varepsilon)\qquad\text{\unsure{i.e}}\\
 \sigma_\Omega\rho^I(\varepsilon)&=\rho^{\sigma I}(\sigma\varepsilon)\sigma_\Omega
 \end{aligned}\right.
\]

Ceci montre que le ss/groupe \inserted{de $\Gamma_0$ (ou $\widetilde\Gamma_0$, ou $\Gamma$, ou $\widetilde\Gamma$)} engendré par $\mathfrak S_3$ et un des $\rho^I(\gamma_0^I)$ (resp. $\rho^I(\gamma^I)$) / \deleted{\unsure{resp}} ou par $\mathfrak S_3$ et \deleted{\unsure{un}} les $3$ $\rho^I(\gamma_0^I)$ (resp $\rho^I(\gamma^I)$) \unsure{sont|est} le $=$, \deleted{\unsure{dits $\Gamma_0$|soit $\Gamma_0$}} $\Gamma_0^+,\widetilde\Gamma_0^+,\Gamma^+,\widetilde\Gamma^+$ ;

% slanted
\margin{/ il suffit $=$ de prendre $\mathfrak S_3$ et \uline{un} des $\rho^I(\gamma_0^{(i)})$ \deleted{$(i\in I)$} (sur les six possibilités)}

et \deleted{\unsure{si on désigne par}} \inserted{le groupe} engendré par les $\rho^I(\gamma_0^I)$ ($\rho^I(\gamma^I)$) en est un ss-groupe invariant \deleted{sont} les $\unsure{\mathbb E_0|\mathrm E_0|\mathbb E\mathrm l_0},\widetilde{\mathbb E}_0,\mathbb E,\widetilde{\mathbb E}$ \deleted{\unsure{$\Gamma_0^+$}} quotient des deux étant \deleted{\unsure{un}} (un groupe quotient de $\mathfrak S_{\unsure{4|3}}$ à priori, mais on voit aisément qu'il est égal à :] $\mathfrak S_3$, en fait c'est un produit $\tfrac12$ direct de $\mathfrak S_3$ et de celui-ci.

% slanted
\margin{(groupe des op. ``élémentaires'')}

\[
 \left\{\begin{aligned}
 \Gamma_0^+&\simeq\mathfrak S_3\underset{1/2}{\cdot}\mathbb E_0&
 \widetilde\Gamma_0^+&\simeq\mathfrak S_3\underset{1/2}{\cdot}\widetilde{\mathbb E}_0\\
 \Gamma^+&\simeq\mathfrak S_3\underset{1/2}{\cdot}\mathbb E&
 \widetilde\Gamma^+&\simeq\mathfrak S_3\underset{1/2}{\cdot}\widetilde{\mathbb E}
 \end{aligned}\right.
\]

\deleted{\uline{\unsure{Questions}}}

\uline{Construction} \unsure{Soient|Soit} $\mathfrak S,\gamma$ deux groupes. Considérons le produit libre $\mathcal L=\mathfrak S\deleted{\unsure{*}}*\gamma$. On a un hom

\[
 \mathcal L=\mathfrak S*\gamma\longrightarrow\mathfrak S
\]

dont le noyau est le ss-groupe \unsure{inv$=$|fermé|inv.} de \inserted{$\mathcal L=$} $\mathfrak S*\gamma$ engendré par $1*\gamma\subset\mathfrak S*\gamma$. Le \unsure{dit} hom admet un scindage \deleted{\unsure{d'un}}\inserted{canonique} $\mathfrak S\hookrightarrow\mathcal L$, d'où
% archive sheet number: 84
